\documentclass[11pt,a4paper]{article}
\usepackage[T1]{fontenc}
\usepackage[margin=27mm]{geometry}
\usepackage[english]{babel}
\usepackage{lmodern,microtype,booktabs,tabularx}
\usepackage{csquotes}
\usepackage[style=apa,backend=biber,natbib=true]{biblatex}
\addbibresource{references.bib}
\usepackage[hidelinks]{hyperref}
\title{Motivations and Psychological Functions of AI-Assisted Music Creation: An Episode-Anchored Survey Protocol}
\author{Mateusz Pomianek}
\date{}
\begin{document}
\maketitle
\begin{abstract}
Generative music systems enable users to produce musical artefacts through prompting, selection and revision, yet the motives preceding these activities need not correspond to their psychological functions or perceived consequences. This prospective cross-sectional study will examine these distinctions in Polish-speaking adult users of generative music systems. Respondents will reconstruct their most recent eligible creative episode within 30 days, regardless of satisfaction or completion. The primary hypothesis concerns the conditional association between the strength of the initial emotional-expression motive and retrospective post-episode emotional clarity, accounting for recalled initial clarity and prespecified covariates. Secondary analyses will investigate mood-oriented motives, exploratory motives, perceived agency, creative operations and adverse experiences. Recruitment will target 480 analytically eligible participants, subject to a fixed recruitment cap and feasibility review before launch. The protocol specifies multiplicity control, missing-data sensitivity analyses and construct boundaries; it makes no causal, therapeutic or psychometric validation claims. No participant data have been collected or analysed.
\end{abstract}
\noindent\textbf{Keywords:} generative AI; music creation; motivation; affect regulation; emotional clarity; survey methodology

\section{Introduction}
Music serves affective, self-related and social functions \parencite{schafer2013}. Research on musical affect regulation distinguishes regulatory goals from strategies and underlying mechanisms \parencite{baltazar2016,baltazar2019}. General emotion regulation theory likewise distinguishes goals, processes and outcomes \parencite{gross2015}. However, evidence concerning listening or conventional music-making cannot be transferred directly to text-conditioned generation. In generative workflows, the user may specify a theme, delegate musical implementation, reject alternatives and listen to a selected output. These operations need not involve the same psychological processes as playing an instrument or passively listening.

Human--AI songwriting research documents the importance of iterative selection and creative collaboration \parencite{huang2020}. Recent research directly addresses task load, psychological ownership and agency in AI collaborative music creation \parencite{he2026}. The present protocol therefore does not claim that creative agency in AI music is an unexplored phenomenon. Its narrower contribution is an episode-anchored distinction between initial motives, reported creative operations, perceived functions and retrospective consequences, including neutral and unsuccessful episodes. The motivational account draws on self-determination theory \parencite{ryan2000}, without treating intrinsic motivation as synonymous with emotional benefit. AIM-ER provides a conceptual vocabulary, not a validated measurement model.

\section{Research questions and hypotheses}
The unit of analysis is one respondent's most recent eligible episode, not a selected favourite song. Initial motives are separately rated and may coexist. Functions refer to what the activity was perceived to enable during the episode; consequences refer to the respondent's perceived state after the episode. Conceptual overlap is reduced by using different prompts and anchors.

\begin{description}
\item[H1 (primary)] Stronger initial motivation to express an emotion is positively associated with recalled post-episode emotional clarity conditional on recalled pre-episode clarity, competing motives and prespecified covariates.
\item[H2 (secondary)] Stronger initial mood-improvement motivation is positively associated with recalled post-episode valence conditional on pre-episode valence and the remaining prespecified covariates.
\item[H3 (secondary)] Stronger initial exploration motivation is positively associated with the number of distinct iterative creative operations performed.
\item[H4 (secondary)] Greater perceived decisional agency is positively associated with perceived personal resonance of the resulting experience, conditional on initial motives and covariates. This is a contemporaneous association, not mediation.
\end{description}
The frequency and distribution of frustration, rumination-like persistence, lack of clarity and negative perceived consequences will be reported descriptively without diagnostic interpretation. The primary hypothesis is not evidence that creating music improves clarity; even the temporal order in the survey is reconstructed retrospectively.

\section{Methods}
\subsection{Design and participants}
A cross-sectional online survey will recruit Polish-speaking adults aged at least 18 who actively prompted, selected, revised or otherwise directed AI-generated music within the past 30 days. Passive listening, non-generative mastering and episodes involving only generic chatbot advice are insufficient. Unfinished and abandoned attempts qualify. Recruitment will use permission-based online communities and university channels, with neutral advertisements that do not promise emotional benefits. This convenience sample does not estimate population prevalence. One response per participant is requested; no IP addresses or identifying platform handles will be retained. Recruitment source will be recorded categorically.

The study is planned as a separate analysis of S02, with a distinct primary outcome and estimand. The same dataset may support P04 or P06 only with transparent dataset identifiers, prespecified non-overlapping primary questions and explicit disclosure; overlapping analyses will be combined rather than presented as independent replications. P05 is reserved for independent measurement validation.

\subsection{Episode and measures}
An episode begins when a respondent intentionally initiates a generative music task and ends when they stop that session, even without an output. The most recent episode is selected before emotional ratings. Respondents may skip items or withdraw before submitting. Initial motive ratings (0--10) cover expression, mood improvement, exploration and entertainment. Distinct operation indicators cover writing a prompt or lyrics, regenerating alternatives, selecting a variant and editing or extending an output. The latter three form an iterative-operation count (0--3); a zero is valid, not missing.

Emotional clarity (0--10; difficulty versus ease understanding one's own feelings) and affective valence (0--10; very unpleasant versus very pleasant) are recalled for just before and just after the episode. Agency concerns perceived influence over creative decisions, while resonance concerns how personally meaningful the experience felt, including experiences without a finished song. Separate descriptive items address perceived expression, exploration, distraction and social connection. Frustration and difficulty disengaging are optional adverse-experience items. These custom single indicators are not clinical measures or validated latent scales. Existing measures of music-related mood regulation \parencite{saarikallio2008} concern related but not identical constructs. No proprietary items will be reproduced or adapted until the Polish version, permissions and scoring instructions are verified.

Covariates are age in years, self-rated musical experience (0--10), typical generator-use days in the last month (0--30), days since the episode (0--30), and recruitment source (university, online community, other). All four motive ratings enter models simultaneously. Platform category and episode completion will be described but not included in the primary adjustment set because they may be downstream of motives. No lyrics, prompts, audio or sensitive narratives are needed.

\subsection{Sample size and stopping}
The smallest association of interest for H1 is a partial correlation of 0.15. A planning approximation based on Fisher's transformation, two-sided alpha .05, power .90 and nine nuisance degrees of freedom gives approximately 473 complete observations; 480 is the operational target. This approximation assumes independent observations, a regular linear model and no severe multicollinearity. It will be checked against a simulation of bounded ratings and missingness before preregistration. Recruitment ends at 480 complete eligible primary records, 560 submissions, or eight weeks after launch, whichever occurs first; only eligibility and completeness counts, never effect estimates, may be monitored. The cap and start date must be frozen before recruitment. Under-recruitment will be disclosed rather than addressed by post-hoc power calculations.

\section{Statistical analysis}
The primary estimand is the population linear-projection coefficient for initial expression motive in the specified sampled-respondent model. The primary analysis uses OLS with HC3 standard errors and a two-sided 95\% confidence interval. The outcome is post-episode clarity; initial clarity, the four motive ratings, age, musical experience, generator-use days, episode recency and recruitment source are included. Recruitment source is represented by two indicators. Report the coefficient in clarity points per one-point motive increase, its interval and partial correlation. H1 uses alpha .05; all other inferential tests are secondary.

H2 substitutes post-valence and initial valence, with the mood motive as the target predictor. H3 uses OLS-HC3 for the bounded operation count (interpreted as a linear mean association); an ordinal model is a sensitivity analysis. H4 uses resonance as the outcome and agency as the target predictor, with four motives and background covariates. Holm adjustment applies across H2--H4. Descriptive adverse items receive denominators and intervals, not a confirmatory hypothesis test. No latent profiles, exploratory factor analysis, structural mediation or causal effects are prespecified.

Report distributions, missingness by item, correlations among motives, variance inflation factors, residual patterns and influential observations. Do not select covariates or exclude observations based on statistical significance. Sensitivity analyses include nonlinear motive effects (restricted spline, descriptive), ordered outcome modelling and excluding episodes older than seven days. The primary complete-record analysis targets respondents with observed required fields; its generalisation depends on the missingness mechanism. If the primary analysis loses any eligible records, multiple imputation under MAR (including outcome, predictors and available auxiliary indicators; 50 imputations) will be compared with the complete-record estimate. A delta sensitivity analysis will shift imputed missing clarity by -2, -1, 0, 1 or 2 points, clipped to bounds. Report uncertainty at each scenario; these assumptions are not empirically identifiable. A missing operation is not coded zero. Data exclusion rules and model failure handling will be frozen before fieldwork.

\section{Ethics, openness and interpretation}
Institutional approval, participant information, informed consent and the applicable data-protection review precede recruitment. Participation is voluntary and may be discontinued without penalty; no intervention, distress induction, platform access or payment is required. Participants may decline questions about emotions. No generated artefacts or personal narratives will be collected. Institutional storage and access rules follow the project data policy. Publish only aggregate results, de-identified code and synthetic example data; row-level release requires a separate disclosure-risk assessment. The protocol, final item wording and analysis code will be versioned and preregistered before recruitment.

Retrospective ratings are subject to reconstructive memory, common-method variance and consistency biases. Conditioning on recalled initial clarity cannot remove these biases, nor can the regression identify a within-person change caused by AI. Selection into generator use, survey participation and episode recall limits external validity. Single-item measures constrain reliability and construct coverage. Neutral, adverse and abandoned episodes will be represented rather than treating perceived benefit as an eligibility criterion.

\section{Results and discussion}
No empirical results are available. The final report will include a participant flow diagram, item-level descriptive statistics, primary and secondary estimates with confidence intervals, missing-data sensitivity results and a transparent account of deviations. Interpretation will distinguish motives from reported psychological functions and outcomes. Results will be evaluated against alternative explanations, including recall bias and prior musical engagement, regardless of the direction or significance of the associations.
\printbibliography
\end{document}
